% Part: normal-modal-logic
% Chapter: tableaux
% Section: soundness

\documentclass[../../../include/open-logic-section]{subfiles}

\begin{document}

\olfileid{nml}{tab}{sou}

\olsection{\Log{K} నిర్దుష్టత}

\begin{editorial}
  ఈ నిర్దుష్టత నిరూపణ సాంప్రదాయ ప్రతిజ్ఞావాక్య తర్కం
  నిర్దుష్టత నిరూపణలోని విషయాన్ని మళ్లీ ఉపయోగిస్తుంది; అంటే
  ప్రతిదానినీ మొదటి నుంచి మళ్లీ నిరూపిస్తుంది. స్వయంపూర్ణ
  నిర్దుష్టత నిరూపణ కావాలంటే ఇది సరే. సాధారణ టాబ్లోల
  నిర్దుష్టతను ఇప్పటికే చూసి ఉంటే ఇది పునరుక్తిగా అనిపిస్తుంది.
  స్వయంపూర్ణ రూపానికీ, మోడల్ కాని సందర్భంపై ఆధారపడే
  రూపానికీ మధ్య ఎంచుకునే వీలు కల్పించాలని యోచిస్తున్నారు.
\end{editorial}

\begin{explain}
  పూర్వసూచిక గల \tetoken{టాబ్లోలు}{!!{tableau}s} నిర్దుష్టమైనవని
  చూపాలంటే, కింది పరికల్పనలకు
  \[
  \sFmla{\True}{!B_1}[1], \dots, \sFmla{\True}{!B_n}[1], \sFmla{\False}{!A}[1]
  \]
  సంవృత \tetoken{టాబ్లో}{!!{tableau}} ఉంటే
  $!B_1, \dots, !B_n \Entails !A$ అని చూపాలి. దాని
  వ్యతిరేకావర్తనాన్ని నిరూపించడం సులువు: ఏదో $\mModel{M}$,
  లోకం~$w$లో అన్ని $i=1$, \dots,~$n$లకు
  $\mSat{M}{!B_i}[w]$ అయినా $\mSat/{M}{!A}[w]$ అయితే,
  ఏ \tetoken{టాబ్లోనూ}{!!{tableau}} సంవృతం చేయలేం.
  \sourcecorrection{OLTENMLTABSOU-001}{ప్రతినమూనాలో~$!A$
    అసత్యం కావాలి; స్థిర మూల వ్యతిరేకావర్తన వాక్యంలో
    సత్య చిహ్నం ఉంది. దాన్ని అసత్య చిహ్నంగా సరిచేశాం.}
  అటువంటి ప్రతినమూనా \tetoken{టాబ్లో}{!!{tableau}}
  ప్రారంభ పరికల్పనలు సంతృప్తిపరచదగినవని చూపుతుంది.
  ఒక \tetoken{టాబ్లో}{!!a{tableau}} శాఖపై పూర్వసూచిక గల
  \tetoken{సూత్రాలన్నీ}{!!{formula}s} సంతృప్తిపరచదగినవి అయితే,
  ఏ నియమాన్ని వర్తింపజేసినా విస్తరించిన శాఖల్లో కనీసం
  ఒకటి అలాగే సంతృప్తిపరచదగినదని చూపడమే నిరూపణ వ్యూహం.
  సంవృత శాఖలు సంతృప్తిపరచలేనివి కాబట్టి,
  సంతృప్తిపరచదగిన పూర్వసూచిక గల \tetoken{సూత్రాల}{!!{formula}s}
  సమితికి చెందిన ఏ \tetoken{టాబ్లోలోనైనా}{!!{tableau}}
  కనీసం ఒక వివృత శాఖ ఉండాలి.

  ఈ వ్యూహాన్ని మోడల్ సందర్భానికి వర్తింపజేయాలంటే,
  ``సంతృప్తిపరచదగినది'' అనే నిర్వచనాన్ని మోడల్ నమూనాలు,
  పూర్వసూచికలకు విస్తరించాలి. అది చేస్తే నిరూపణ సూటిగా సాగుతుంది.
\end{explain}

\begin{defn}
  $P$ ఏదో ఒక పూర్వసూచికల సమితి, అంటే $P \subseteq (\PosInt)^*
  \setminus \{\emptyseq\}$; $\mModel{M}$ ఒక నమూనా అనుకుందాం.
  $\sigma$, $\sigma.n$ రెండూ~$P$లో ఉన్నప్పుడల్లా
  $Rf(\sigma)f(\sigma.n)$ అయితే, ప్రమేయకం~$f\colon P \to W$ను
  $\mModel{M}$లో~$P$కు \emph{అర్థనిర్దేశం} అంటాం.

  పూర్వసూచికల~$P$ అర్థనిర్దేశాన్ని బట్టి ఇలా నిర్వచించవచ్చు:
  \begin{enumerate}
  \item $\mSat{M}{!A}[f(\sigma)]$ అయితే, అప్పుడు మాత్రమే
    $\mModel{M}$, $\sFmla{\True}{!A}[\sigma]$ను సంతృప్తిపరుస్తుంది.
  \item $\mSat/{M}{!A}[f(\sigma)]$ అయితే, అప్పుడు మాత్రమే
    $\mModel{M}$, $\sFmla{\False}{!A}[\sigma]$ను సంతృప్తిపరుస్తుంది.
  \end{enumerate}
\end{defn}

\begin{defn}
  $\Gamma$ పూర్వసూచిక గల \tetoken{సూత్రాల}{!!{formula}s}
  సమితి అనుకుందాం; అందులో కనిపించే పూర్వసూచికల సమితి
  $P(\Gamma)$ అనుకుందాం. $f$ అనేది $\mModel{M}$లో
  $P(\Gamma)$కు అర్థనిర్దేశమైతే, $\Gamma$లోని ప్రతి
  పూర్వసూచిక గల \tetoken{సూత్రాన్ని}{!!{formula}}~$\mModel{M}$,
  $f$ను బట్టి సంతృప్తిపరిస్తే, $\mModel{M}$, $f$ను బట్టి
  $\Gamma$ను సంతృప్తిపరుస్తుంది, $\mSat{M}{\Gamma}[f]$,
  అంటాం. $\mSat{M}{\Gamma}[f]$ అయ్యే నమూనా~$\mModel{M}$,
  $P(\Gamma)$ అర్థనిర్దేశం~$f$ ఉంటే, అప్పుడు మాత్రమే
  $\Gamma$ \emph{సంతృప్తిపరచదగినది}.
\end{defn}

\begin{prop}
  ఏదో \tetoken{సూత్రం}{!!{formula}}~$!A$,
  పూర్వసూచిక~$\sigma$ కోసం $\Gamma$లో
  $\sFmla{\True}{!A}[\sigma]$, $\sFmla{\False}{!A}[\sigma]$
  రెండూ ఉంటే $\Gamma$ సంతృప్తిపరచలేనిది.
\end{prop}

\begin{proof}
  $\mSat{M}{!A}[f(\sigma)]$,
  $\mSat/{M}{!A}[f(\sigma)]$ రెండూ అయ్యే నమూనా~$\mModel{M}$,
  $P(\Gamma)$ అర్థనిర్దేశం~$f$ ఉండవు.
\end{proof}

\begin{thm}[నిర్దుష్టత]
  \ollabel{thm:tableau-soundness}
  $\Gamma$కు సంవృత \tetoken{టాబ్లో}{!!{tableau}} ఉంటే,
  $\Gamma$ సంతృప్తిపరచలేనిది.
\end{thm}

\begin{proof}
ఒక \tetoken{టాబ్లో}{!!a{tableau}} శాఖపై ఉన్న
\tetoken{చిహ్నిత సూత్రాల}{!!{signed formula}s} సమితి
సంతృప్తిపరచదగినదైతే, అప్పుడు మాత్రమే ఆ శాఖను
సంతృప్తిపరచదగినది అంటాం; కనీసం ఒక అటువంటి శాఖ ఉన్న
\tetoken{టాబ్లోను}{!!a{tableau}} సంతృప్తిపరచదగినది అంటాం.

మనం చూపేది ఇది: సంతృప్తిపరచదగిన \tetoken{టాబ్లోను}{!!{tableau}}
నిగమన నియమాల్లో ఒకదానితో విస్తరించినప్పుడు వచ్చే
\tetoken{టాబ్లో}{!!{tableau}} కూడా సంతృప్తిపరచదగినదే.
దీనితో సిద్ధాంతం నిరూపితమవుతుంది: ప్రతి సంవృత
\tetoken{టాబ్లో}{!!{tableau}},~$\Gamma$లోని పరికల్పనలు మాత్రమే
గల \tetoken{టాబ్లోకు}{!!{tableau}} నిగమన నియమాలు వర్తింపజేయడం
వల్ల ఏర్పడుతుంది. కనుక $\Gamma$ సంతృప్తిపరచదగినదైతే
దాని కోసం ప్రతి \tetoken{టాబ్లో}{!!{tableau}} కూడా
సంతృప్తిపరచదగినదే. కానీ సంవృత \tetoken{టాబ్లో}{!!{tableau}}
స్పష్టంగా సంతృప్తిపరచదగినది కాదు: దాని శాఖలన్నీ
సంవృతం, సంవృత శాఖలు సంతృప్తిపరచలేనివి.

మనకు సంతృప్తిపరచదగిన \tetoken{టాబ్లో}{!!{tableau}}, అంటే
కనీసం ఒక సంతృప్తిపరచదగిన శాఖ ఉన్న \tetoken{టాబ్లో}{!!a{tableau}}
ఉందనుకుందాం. నిగమన నియమం వర్తింపజేస్తే ఒక శాఖకు
\tetoken{చిహ్నిత సూత్రాలు}{!!{signed formula}s} చేర్చబడతాయి
లేదా శాఖ రెండుగా విడిపోతుంది. నియమం వల్ల విస్తరించని
సంతృప్తిపరచదగిన శాఖ \tetoken{టాబ్లోలో}{!!{tableau}} ఉంటే,
విస్తరించిన \tetoken{టాబ్లోలోనూ}{!!{tableau}} అది అలాగే
ఉంటుంది; కాబట్టి విస్తరించిన టాబ్లో సంతృప్తిపరచదగినది.
అందువల్ల సంతృప్తిపరచదగిన శాఖకే నియమాన్ని వర్తింపజేసే
సందర్భాన్ని మాత్రమే పరిశీలించాలి.

ఆ శాఖపైని \tetoken{చిహ్నిత సూత్రాల}{!!{signed formula}s}
సమితి~$\Gamma$ అనుకుందాం; నియమాన్ని వర్తింపజేసే
\tetoken{చిహ్నిత సూత్రం}{!!{signed formula}}~$\sFmla{S}{!A}[\sigma] \in \Gamma$
అనుకుందాం. నియమం వల్ల శాఖ విడిపోకపోతే, విస్తరించిన
శాఖ, అంటే నియమ నిష్కర్షలతో కలిసిన~$\Gamma$, ఇంకా
సంతృప్తిపరచదగినదని చూపాలి. శాఖ విడిపోతే, ఫలితంగా
వచ్చే రెండు శాఖల్లో కనీసం ఒకటి సంతృప్తిపరచదగినదని
చూపాలి.
\tagfalse{prvDiamond}
మొదట ఒకే పూర్వాధారం గల నిగమనాలను పరిశీలిద్దాం.
\begin{enumerate}
\item $\sFmla{\True}{\lnot !B}[\sigma] \in \Gamma$కు
  $\TRule{\True}{\lnot}$ వర్తింపజేసి శాఖను విస్తరిస్తాం.
  అప్పుడు విస్తరించిన శాఖలోని \tetoken{చిహ్నిత సూత్రాల}{!!{signed formula}s}
  సమితి $\Gamma \cup \{\sFmla{\False}{!B}[\sigma]\}$.
  $\mSat{M}{\Gamma}[f]$ అనుకుందాం. ముఖ్యంగా,
  $\mSat{M}{\lnot !B}[f(\sigma)]$. అందువల్ల
  $\mSat/{M}{!B}[f(\sigma)]$; అంటే~$f$ను బట్టి $\mModel{M}$,
  $\sFmla{\False}{!B}[\sigma]$ను సంతృప్తిపరుస్తుంది.
\item $\sFmla{\False}{\lnot !B}[\sigma] \in \Gamma$కు
  $\TRule{\False}{\lnot}$ వర్తింపజేసి శాఖను విస్తరిస్తాం:
  వ్యాయామం.
\item $\sFmla{\True}{!B \land !C}[\sigma] \in \Gamma$కు
  $\TRule{\True}{\land}$ వర్తింపజేస్తే శాఖపై రెండు కొత్త
  \tetoken{చిహ్నిత సూత్రాలు}{!!{signed formula}s},
  $\sFmla{\True}{!B}[\sigma]$, $\sFmla{\True}{!C}[\sigma]$ వస్తాయి.
  $\mSat{M}{\Gamma}[f]$, ముఖ్యంగా $\mSat{M}{!B \land
    !C}[f(\sigma)]$ అనుకుందాం. అప్పుడు $\mSat{M}{!B}[f(\sigma)]$,
  $\mSat{M}{!C}[f(\sigma)]$. అందువల్ల~$f$ను బట్టి $\mModel{M}$,
  $\sFmla{\True}{!B}[\sigma]$, $\sFmla{\True}{!C}[\sigma]$
  రెండింటినీ సంతృప్తిపరుస్తుంది.
\item $\sFmla{\False}{!B \lor !C}[\sigma] \in \Gamma$కు
  $\TRule{\False}{\lor}$ వర్తింపజేసి శాఖను విస్తరిస్తాం:
  వ్యాయామం.
  \sourcecorrection{OLTENMLTABSOU-002}{ఈ అంశంలో మూలం
    $\sFmla{\False}{!B \lor !C}$కు పూర్వసూచికను వదిలింది.
    ఇదే అంశాల వరుసకు, నియమ నిర్వచనానికి అనుగుణంగా
    $[\sigma]$ను పునరుద్ధరించాం.}
\item $\sFmla{\False}{!B \lif !C}[\sigma] \in \Gamma$కు
  $\TRule{\False}{\lif}$ వర్తింపజేస్తాం. అప్పుడు శాఖపై
  రెండు కొత్త \tetoken{చిహ్నిత సూత్రాలు}{!!{signed formula}s}:
  $\sFmla{\True}{!B}[\sigma]$, $\sFmla{\False}{!C}[\sigma]$.
  $\mSat{M}{\Gamma}[f]$, ముఖ్యంగా
  $\mSat/{M}{!B \lif !C}[f(\sigma)]$ అనుకుందాం. అప్పుడు
  $\mSat{M}{!B}[f(\sigma)]$, $\mSat/{M}{!C}[f(\sigma)]$.
  కాబట్టి~$\mModel{M}, f$, $\sFmla{\True}{!B}[\sigma]$,
  $\sFmla{\False}{!C}[\sigma]$ రెండింటినీ సంతృప్తిపరుస్తుంది.

\iftag{prvBox}{%
\item $\sFmla{\True}{\Box !B}[\sigma] \in \Gamma$కు
  $\TRule{\True}{\Box}$ వర్తింపజేసి శాఖను విస్తరిస్తాం:
  \iftag{probBox}{వ్యాయామం.}{దీనివల్ల శాఖపై కొత్త
    \tetoken{చిహ్నిత సూత్రం}{!!{signed
      formula}}~$\sFmla{\True}{!B}[\sigma.n]$ వస్తుంది;
    ఇక్కడ $\sigma.n \in P(\Gamma)$ (ఎందుకంటే $\sigma.n$
    ఉపయోగించిన పూర్వసూచిక కావాలి). $\mSat{M}{\Gamma}[f]$,
    ముఖ్యంగా $\mSat{M}{\Box !B}[f(\sigma)]$ అనుకుందాం.
    $f$ పూర్వసూచికల అర్థనిర్దేశం, $\sigma$,
    $\sigma.n \in P(\Gamma)$ కాబట్టి $Rf(\sigma)f(\sigma.n)$.
    అందువల్ల $\mSat{M}{!B}[f(\sigma.n)]$; అంటే
    $\mModel{M}, f$, $\sFmla{\True}{!B}[\sigma.n]$ను
    సంతృప్తిపరుస్తుంది.}

\item $\sFmla{\False}{\Box !B}[\sigma] \in \Gamma$కు
  $\TRule{\False}{\Box}$ వర్తింపజేసి శాఖను విస్తరిస్తాం:
  \iftag{probBox}{వ్యాయామం.}{దీనివల్ల కొత్త
    \tetoken{చిహ్నిత సూత్రం}{!!{signed
      formula}}~$\sFmla{\False}{!B}[\sigma.n]$ వస్తుంది;
    ఇక్కడ $\sigma.n$ శాఖపై కొత్త పూర్వసూచిక, అంటే
    $\sigma.n \notin P(\Gamma)$.
    \sourcecorrection{OLTENMLTABSOU-003}{ఈ $\TRule{\False}{\Box}$ అంశంలో
      మూలం కొత్త నిష్కర్షను $\sFmla{\False}{!A}[\sigma.n]$గా
      వ్రాసింది; పరికల్పన $\Box !B$, మిగతా నిరూపణకు అనుగుణంగా
      $\sFmla{\False}{!B}[\sigma.n]$గా సరిచేశాం.
      ఇదే అంశంలోని నిర్మాణ/సంతృప్తి చిహ్నాలను మోడల్
      సందర్భానికనుగుణంగా ఏకరీతిగా వాడాం.}
    $\Gamma$ సంతృప్తిపరచదగినది కాబట్టి,
    $\mSat{M}{\Gamma}[f]$ అయ్యే నమూనా~$\mModel{M}$,
    $P(\Gamma)$ అర్థనిర్దేశం~$f$ ఉన్నాయి; ముఖ్యంగా
    $\mSat/{M}{\Box !B}[f(\sigma)]$. $\Gamma \cup
    \{\sFmla{\False}{!B}[\sigma.n]\}$ సంతృప్తిపరచదగినదని
    చూపాలి. ఇందుకోసం $P(\Gamma) \cup \{\sigma.n\}$కు
    అర్థనిర్దేశాన్ని ఇలా నిర్వచిద్దాం:

  $\mSat/{M}{\Box !B}[f(\sigma)]$ కాబట్టి, $Rf(\sigma)w$,
  $\mSat/{M}{!B}[w]$ అయ్యే $w \in W$ ఉంది. $f'(\sigma.n) = w$
  తప్ప మిగతా చోట్ల~$f$లాగే ఉండే~$f'$ను తీసుకుందాం.
  $f'(\sigma) = f(\sigma)$, $Rf(\sigma)w$ కాబట్టి
  $Rf'(\sigma)f'(\sigma.n)$; కనుక~$f'$ అనేది
  $P(\Gamma) \cup \{\sigma.n\}$కు అర్థనిర్దేశం. స్పష్టంగా
  $\mSat/{M}{!B}[f'(\sigma.n)]$. అన్ని
  $\sigma' \in P(\Gamma)$లకు $f(\sigma') = f'(\sigma')$ కాబట్టి
  $\mSat{M}{\Gamma}[f']$. అందువల్ల~$\mModel{M}, f'$,
  $\Gamma \cup \{\sFmla{\False}{!B}[\sigma.n]\}$ను
  సంతృప్తిపరుస్తుంది.}
  }{}

\iftag{prvDiamond}{%
\item $\sFmla{\True}{\Diamond !B}[\sigma] \in \Gamma$కు
  $\TRule{\True}{\Diamond}$ వర్తింపజేసి శాఖను విస్తరిస్తాం:
  \iftag{probDiamond}{వ్యాయామం.}{దీనివల్ల కొత్త
    \tetoken{చిహ్నిత సూత్రం}{!!{signed
      formula}}~$\sFmla{\True}{!B}[\sigma.n]$ వస్తుంది;
    $\sigma.n$ శాఖపై కొత్త పూర్వసూచిక, అంటే
    $\sigma.n \notin P(\Gamma)$.
    \sourcecorrection{OLTENMLTABSOU-004}{ఈ $\TRule{\True}{\Diamond}$ అంశంలో
      మూలం కొత్త నిష్కర్షను $\sFmla{\True}{!A}[\sigma.n]$గా
      వ్రాసింది; పరికల్పన $\Diamond !B$, మిగతా నిరూపణకు అనుగుణంగా
      $\sFmla{\True}{!B}[\sigma.n]$గా సరిచేశాం.
      ఇదే అంశంలోని నిర్మాణ/సంతృప్తి చిహ్నాలను మోడల్
      సందర్భానికనుగుణంగా ఏకరీతిగా వాడాం.}
    $\Gamma$ సంతృప్తిపరచదగినది కాబట్టి,
    $\mSat{M}{\Gamma}[f]$ అయ్యే నమూనా~$\mModel{M}$,
    $P(\Gamma)$ అర్థనిర్దేశం~$f$ ఉన్నాయి; ముఖ్యంగా
    $\mSat{M}{\Diamond !B}[f(\sigma)]$. $\Gamma \cup
    \{\sFmla{\True}{!B}[\sigma.n]\}$ సంతృప్తిపరచదగినదని
    చూపాలి. ఇందుకోసం $P(\Gamma) \cup \{\sigma.n\}$కు
    అర్థనిర్దేశాన్ని ఇలా నిర్వచిద్దాం:

  $\mSat{M}{\Diamond !B}[f(\sigma)]$ కాబట్టి, $Rf(\sigma)w$,
  $\mSat{M}{!B}[w]$ అయ్యే $w \in W$ ఉంది. $f'(\sigma.n) = w$
  తప్ప మిగతా చోట్ల~$f$లాగే ఉండే~$f'$ను తీసుకుందాం.
  $f'(\sigma) = f(\sigma)$, $Rf(\sigma)w$ కాబట్టి
  $Rf'(\sigma)f'(\sigma.n)$; కనుక~$f'$ అనేది
  $P(\Gamma) \cup \{\sigma.n\}$కు అర్థనిర్దేశం. స్పష్టంగా
  $\mSat{M}{!B}[f'(\sigma.n)]$. అన్ని
  $\sigma' \in P(\Gamma)$లకు $f(\sigma') = f'(\sigma')$ కాబట్టి
  $\mSat{M}{\Gamma}[f']$. అందువల్ల~$\mModel{M}, f'$,
  $\Gamma \cup \{\sFmla{\True}{!B}[\sigma.n]\}$ను
  సంతృప్తిపరుస్తుంది.}
\item $\sFmla{\False}{\Diamond !B}[\sigma] \in \Gamma$కు
  $\TRule{\False}{\Diamond}$ వర్తింపజేసి శాఖను విస్తరిస్తాం:
  \iftag{probDiamond}{వ్యాయామం.}{దీనివల్ల శాఖపై కొత్త
    \tetoken{చిహ్నిత సూత్రం}{!!{signed
      formula}}~$\sFmla{\False}{!B}[\sigma.n]$ వస్తుంది;
    ఇక్కడ $\sigma.n \in P(\Gamma)$ (ఎందుకంటే $\sigma.n$
    ఉపయోగించిన పూర్వసూచిక కావాలి). $\mSat{M}{\Gamma}[f]$,
    ముఖ్యంగా $\mSat/{M}{\Diamond !B}[f(\sigma)]$ అనుకుందాం.
    $f$ పూర్వసూచికల అర్థనిర్దేశం, $\sigma$,
    $\sigma.n \in P(\Gamma)$ కాబట్టి $Rf(\sigma)f(\sigma.n)$.
    అందువల్ల $\mSat/{M}{!B}[f(\sigma.n)]$; అంటే
    $\mModel{M}, f$, $\sFmla{\False}{!B}[\sigma.n]$ను
    సంతృప్తిపరుస్తుంది.}
  }{}
\end{enumerate}
ఇప్పుడు శాఖను రెండుగా విడగొట్టే నిగమనాలను పరిశీలిద్దాం.
\sourcecorrection{OLTENMLTABSOU-005}{ఇక్కడి నియమాలకు రెండు
  పూర్వాధారాలు కాదు, ఒక్క పూర్వాధారం నుంచి రెండు శాఖలు
  వస్తాయి. మూలంలోని ``two premises''ను ఆ శాఖల
  అర్థానికి అనుగుణంగా సరిచేశాం.}
\begin{enumerate}
\item $\sFmla{\False}{!B \land !C}[\sigma] \in \Gamma$కు
  $\TRule{\False}{\land}$ వర్తింపజేస్తే శాఖ రెండుగా విడిపోతుంది:
  ఎడమ శాఖ $\sFmla{\False}{!B}[\sigma]$ ద్వారా, కుడి శాఖ
  $\sFmla{\False}{!C}[\sigma]$ ద్వారా కొనసాగుతాయి.
  $\mSat{M}{\Gamma}[f]$, ముఖ్యంగా
  $\mSat/{M}{!B \land !C}[f(\sigma)]$ అనుకుందాం. అప్పుడు
  $\mSat/{M}{!B}[f(\sigma)]$ లేదా $\mSat/{M}{!C}[f(\sigma)]$.
  మొదటి సందర్భంలో~$\mModel{M}, f$,
  $\sFmla{\False}{!B}[\sigma]$ను సంతృప్తిపరుస్తుంది;
  అంటే ఎడమ శాఖ సంతృప్తిపరచదగినది. రెండో సందర్భంలో
  $\mModel{M}, f$, $\sFmla{\False}{!C}[\sigma]$ను
  సంతృప్తిపరుస్తుంది; అంటే కుడి శాఖ సంతృప్తిపరచదగినది.
\item $\sFmla{\True}{!B \lor !C}[\sigma] \in \Gamma$కు
  $\TRule{\True}{\lor}$ వర్తింపజేసి శాఖను విస్తరిస్తాం:
  వ్యాయామం.
\item $\sFmla{\True}{!B \lif !C}[\sigma] \in \Gamma$కు
  $\TRule{\True}{\lif}$ వర్తింపజేసి శాఖను విస్తరిస్తాం:
  వ్యాయామం.
\end{enumerate}
\end{proof}

\begin{prob}
\olref[nml][tab][sou]{thm:tableau-soundness} నిరూపణను పూర్తి చేయండి.
\end{prob}

\begin{cor}
\ollabel{cor:entailment-soundness}
$\Gamma \Proves !A$ అయితే $\Gamma \Entails !A$.
\end{cor}

\begin{proof}
  $\Gamma \Proves !A$ అయితే, ఏవో $!B_1$, \dots,
  $!B_n \in \Gamma$ కోసం $\Delta =
  \{\sFmla{\False}{!A}[1], \sFmla{\True}{!B_1}[1],
  \dots, \sFmla{\True}{!B_n}[1]\}$కు సంవృత
  \tetoken{టాబ్లో}{!!{tableau}} ఉంది. $\Gamma \Entails !A$
  అని చూపాలి. అలా కాదనుకుందాం. అప్పుడు ఏదో
  $\mModel{M}$, $w$ కోసం అన్ని $i=1$, \dots,~$n$లకు
  $\mSat{M}{!B_i}[w]$, కానీ $\mSat/{M}{!A}[w]$.
  $f(1) = w$గా తీసుకుందాం; అప్పుడు~$f$ అనేది
  $\mModel{M}$లో~$P(\Delta)$కు అర్థనిర్దేశం; $f$ను బట్టి
  $\mModel{M}$,~$\Delta$ను సంతృప్తిపరుస్తుంది. కానీ
  \olref{thm:tableau-soundness} ప్రకారం సంవృత
  \tetoken{టాబ్లో}{!!{tableau}} ఉన్న~$\Delta$
  సంతృప్తిపరచలేనిది; ఇది వైరుధ్యం. కనుక
  $\Gamma \Entails !A$ తప్పనిసరి.
  \sourcecorrection{OLTENMLTABSOU-006}{ఈ నిరూపణ చివర మూలం
    పరికల్పన $\Gamma \Proves !A$నే మళ్లీ నిష్కర్షగా
    వ్రాసింది. వైరుధ్య నిరూపణ చూపింది, సిద్ధాంతం చెప్పేది
    $\Gamma \Entails !A$; దానిని పునరుద్ధరించాం.}
\end{proof}

\begin{cor}
\ollabel{cor:weak-soundness}
$\Proves !A$ అయితే అన్ని నమూనాల్లో~$!A$ సత్యం.
\end{cor}

\end{document}
